\documentclass{article}
% \documentclass[twocolumn]{article}
\usepackage{graphicx} % Required for inserting images
\usepackage{tikz} 
\usepackage{amsmath}
\usepackage{booktabs}
\usepackage{amssymb}
\usepackage{amsthm}
\usepackage{mathtools}

\newtheorem{theorem}{Theorem}[section]

\newcommand{\mybox}[5]{
    \begin{center}
    \begin{tikzpicture}
        \node[anchor=text,text width=\columnwidth+1.0cm, draw, rounded corners, line width=1pt, fill=#3, inner sep=5mm] (big) {\\#4};
        \node[draw, rounded corners, line width=.5pt, fill=#2, anchor=west, xshift=5mm] (small) at (big.north west) {#1};
    \end{tikzpicture}
    % \caption{#5}
    \end{center}
}


\title{Mithril in Game-based Syntax}

\author{Wuyunsiqin \quad \quad Bingsheng Zhang \\ 
Zhejiang University \\
oyun@zju.edu.cn \quad bingsheng@zju.edu.cn}
% \date{May 20 2025}




\begin{document}

\maketitle



\section{Preliminaries}
\subsection{Notations}

There are some basic notation and conventions used throughout this work.

\begin{itemize}
  \item The security parameter is denoted by $\lambda \in \mathbb{N}$.
  %  All probabilistic polynomial-time (PPT) algorithms and adversaries are assumed to run in time polynomial in $\lambda$. Whenever an algorithm takes $1^\lambda$ as input, this denotes the unary encoding of $\lambda$.

  % \item For an integer $n \in \mathbb{N}$, we write $[n] := \{1,\dots,n\}$ for the index set of $n$ parties. If $\mathbf{x} = (x_1,\dots,x_n)$ is a vector indexed by $[n]$, we use boldface letters for vectors (for example, $\mathbf{vk}$ and $\mathbf{w}$) and plain letters for their components.

  \item For any finite set $S$, we write $|S|$ for its cardinality. We use the notation
        \[
          x \xleftarrow{\$} S
        \]
        to indicate that $x$ is sampled uniformly at random from the finite set $S$.

  % \item For an integer $q \ge 2$, we write $\mathbb{Z}_q$ for the ring of integers modulo~$q$ and $\mathbb{Z}_q^\ast$ for its multiplicative group of invertible elements.

  % \item If $X$ denotes a random variable or a probabilistic experiment, we write $\Pr[X \in E]$ for the probability of an event $E$. When the underlying experiment is clear from the context, we abbreviate this to $\Pr[\cdot]$.

  \item A function $\mathsf{negl} : \mathbb{N} \to \mathbb{R}_{\ge 0}$ is called \emph{negligible} if for every polynomial $p(\cdot)$ there exists $N \in \mathbb{N}$ such that for all $\lambda \ge N$ it holds that $\mathsf{negl}(\lambda) < 1/p(\lambda)$. We write $\mathsf{negl}(\lambda)$ to denote an unspecified negligible function in the security parameter~$\lambda$.
\end{itemize}


\subsection{BLS Signature Scheme}

Let $\lambda$ be the security parameter. We assume a bilinear group
generation algorithm which, on input $1^\lambda$, outputs
\[
  (q, G_1, G_2, G_T, g_1, g_2, e),
\]
where $G_1,G_2,G_T$ are cyclic groups of prime order $q$, $g_1\in G_1$,
$g_2\in G_2$ are generators, and
$e : G_1 \times G_2 \to G_T$ is a bilinear pairing. Let
$H : \{0,1\}^* \to G_1$ be a hash function modeled as a random oracle.

We write the BLS signature scheme as
\[
  \Pi_{\mathsf{BLS}} = (\mathsf{Setup}, \mathsf{KGen}, \mathsf{Sign}, \mathsf{Ver}).
\]

\begin{itemize}
  \item $pp \leftarrow \mathsf{Setup}(1^\lambda)$:  
        run the group generator to obtain
        $(q, G_1, G_2, G_T, g_1, g_2, e)$ and fix $H$.
        Output public parameters
        \[
          \mathsf{pp} = (q, G_1, G_2, G_T, g_1, g_2, e, H).
        \]

  \item $(sk,vk) \leftarrow \mathsf{KGen}(\mathsf{pp})$:  
        sample a secret key $sk \xleftarrow{\$} \mathbb{Z}_q$, where
        $\mathbb{Z}_q$ denotes the integers modulo $q$, and set
        \[
          vk = g_2^{sk} \in G_2.
        \]
        Output the key pair $(sk, vk)$.

  \item $\sigma \leftarrow \mathsf{Sign}(\mathsf{pp}, sk, msg)$:  
        on input a message $msg \in \{0,1\}^*$, compute
        $h \leftarrow H(msg) \in G_1$ and set
        \[
          \sigma = h^{sk} \in G_1.
        \]
        Output the signature $\sigma$.

  \item $b \leftarrow \mathsf{Ver}(\mathsf{pp}, vk, msg, \sigma)$, where $b \in \{0,1\}$:  
        compute $h \leftarrow H(msg)$ and output
        \[
          b = 1
          \;\text{iff}\;
          e(\sigma, g_2) = e(h, vk).
        \]
\end{itemize}

\subsection{Merkle Tree}

Let $H_{\mathsf{MT}}  : \{0,1\}^* \to \{0,1\}^\lambda$ be a collision-resistant hash function.

We denote the Merkle tree primitive by
\[
  \Pi_{\mathsf{MT}} = (\mathsf{Create}, \mathsf{Check}).
\]

\begin{itemize}

  \item $\mathsf{avk} \leftarrow \mathsf{Create}(\mathbf{vk}, \mathbf{w})$:
        on input a vector of verification keys
        $\mathbf{vk} = (vk_1,\dots,vk_N)$ and a matching stake vector
        $\mathbf{w}=(w_1,\dots,w_N)$, where $N \in \mathbb{N}$ denotes
        the number of parties, set leaves
        $\ell_i := H_{\mathsf{MT}}(vk_i , w_i)$ and output the Merkle root
        $\mathsf{avk}$, which we also call an aggregate verification key.

  \item $b \leftarrow \mathsf{Check}(\mathsf{avk}, vk, w, path)$, where $b \in \{0,1\}$:
        on input a root $\mathsf{avk}$, a pair $(vk,w)$, and an authentication
        path $path$, recompute the root from the leaf
        $\ell := H_{\mathsf{MT}}(vk , w)$ and $path$. Output $b = 1$ iff the
        recomputed root equals $\mathsf{avk}$, and $b = 0$ otherwise.

\end{itemize}


\section{Stake-based threshold multi-signature}


\subsection{Syntax}

We consider a stake-based threshold multi-signature (STMS) scheme whose acceptance is determined on a per-stake basis via a ``$t$ wins in $n$ lotteries'' parametrization, for some integers $t,n \in \mathbb{N}$.



\[
  \Pi_{\mathsf{STMS}}
  =
  (\mathsf{Setup}, \mathsf{KGen}, \mathsf{AKey}, \mathsf{ACheck},
   \mathsf{Sig}, \mathsf{Ver}, \mathsf{ASig}, \mathsf{AVer}) .
\]

\begin{itemize}
  \item $\mathsf{pp} \leftarrow \mathsf{Setup}(1^\lambda)$:  
        on input a security parameter $1^\lambda$, outputs public system parameters $\mathsf{pp}$.  
        These parameters are implicit inputs to all other algorithms.

  \item $(sk,\, vk) \leftarrow \mathsf{KGen}(\mathsf{pp})$:  
        on input the public parameters $\mathsf{pp}$, generates a signing/verification key pair $(sk,vk)$ for a single party.

  \item $\mathsf{avk} \leftarrow \mathsf{AKey}(\mathbf{vk},\, \mathbf{w})$:  
        given a list of verification keys $\mathbf{vk} := (vk_1,\ldots,vk_N)$ and a matching stake vector $\mathbf{w} := (w_1,\ldots,w_N)$,  
        where $N \in \mathbb{N}$ denotes the number of parties, deterministically produces an aggregate verification key $\mathsf{avk}$.

  \item $b \leftarrow \mathsf{ACheck}(\mathsf{avk},\, \mathbf{vk},\, \mathbf{w})$, where $b \in \{0,1\}$:  
        deterministically checks whether $\mathsf{avk}$ is a valid aggregate verification key for $(\mathbf{vk},\mathbf{w})$,  
        outputting $b = 1$ if it accepts and $b = 0$ otherwise.

  \item $\sigma \leftarrow \mathsf{Sig}(vk,\, w,\, sk,\, j,\, m)$:  
        attempts to produce a signature $\sigma$ for message $m$ and lottery index $j \in [n]$  
        under verification key $vk$, secret key $sk$, and stake $w$.  
        It outputs $\bot$ if the user is ineligible for $(m,j)$ under stake $w$.

  \item $b \leftarrow \mathsf{Ver}(vk,\, w,\, j,\, m,\, \sigma)$, where $b \in \{0,1\}$:  
        given a verification key $vk$, a stake value $w$, a lottery index $j$, a message $m$, and a signature $\sigma$,  
        deterministically checks whether $\sigma$ is a valid signature for $(m,j)$ under $(vk,w)$  
        and outputs $b = 1$ if it accepts and $b = 0$ otherwise.

  \item $\tau \leftarrow \mathsf{ASig}(\mathbf{vk},\, \mathbf{w},\, m,\, \{(vk_i,\sigma_i)\}_{i=1}^{t})$:  
        given $(\mathbf{vk},\mathbf{w})$, a message $m$, and a multiset $\{(vk_i,\sigma_i)\}_{i=1}^{t}$ of signatures,  
        attempts to produce an aggregate signature $\tau$; it outputs $\bot$ if the provided inputs are insufficient or inconsistent.

  \item $b \leftarrow \mathsf{AVer}(\mathsf{avk},\, m,\, \tau,\, T)$, where $b \in \{0,1\}$:  
        given a message $m$, an aggregate verification key $\mathsf{avk}$, an aggregate signature $\tau$,  
        and a stake threshold parameter $T$, outputs $b = 1$ if $\tau$ is accepted as a valid aggregate  
        (stake-based threshold) multi-signature on $m$ under $\mathsf{avk}$ and threshold $T$, and $b = 0$ otherwise.

\end{itemize}



\subsubsection{Correctness}

Let 
\[
  \Pi_{\mathsf{STMS}}
  =
  (\mathsf{Setup}, \mathsf{KGen}, \mathsf{AKey}, \mathsf{ACheck},
   \mathsf{Sig}, \mathsf{Ver}, \mathsf{ASig}, \mathsf{AVer})
\]
be a stake-based threshold multi-signature (STMS) scheme as specified above.  
Let $\mathcal{M}$ denote the message space.

\paragraph{Definition (Correctness).}
We say that $\Pi_{\mathsf{STMS}}$ is \emph{correct} if for every security parameter $\lambda \in \mathbb{N}$,
every number of parties $N \in \mathbb{N}$, every vector of verification keys
$\mathbf{vk} = (vk_1,\dots,vk_N)$ generated honestly by $\mathsf{KGen}$, every stake vector
$\mathbf{w} = (w_1,\dots,w_N)$, every message $m \in \mathcal{M}$, every stake threshold parameter
$T$, and every lottery index $j$ in the relevant index set, the following properties hold except with
negligible probability in $\lambda$ over the randomness of all algorithms involved.


\begin{itemize}
  \item \textbf{Correctness of aggregate key.}
        If $\mathsf{avk} \leftarrow \mathsf{AKey}(\mathbf{vk}, \mathbf{w})$, then
        \[
          \mathsf{ACheck}(\mathsf{avk}, \mathbf{vk}, \mathbf{w}) = 1.
        \]

  \item \textbf{Correctness of individual signatures.}
        For every party index $i \in [N]$, let $(sk_i, vk_i)$ be generated by
        $(sk_i, vk_i) \leftarrow \mathsf{KGen}(\mathsf{pp})$.
        If $\sigma_i \leftarrow \mathsf{Sig}(vk_i, w_i, sk_i, j, m)$ and $\sigma_i \neq \bot$, then
        \[
          \mathsf{Ver}(vk_i, w_i, j, m, \sigma_i) = 1.
        \]

  \item \textbf{Correctness of aggregate signatures.}
        Let $t \in \mathbb{N}$ and, for each $i \in [t]$, let $(sk_i, vk_i)$ be generated by
        $(sk_i, vk_i) \leftarrow \mathsf{KGen}(\mathsf{pp})$ and let $w_i$ be the corresponding stake value.
        For a message $m$ and a lottery index $j$, let
        $\sigma_i \leftarrow \mathsf{Sig}(vk_i, w_i, sk_i, j, m)$ be such that $\sigma_i \neq \bot$ for all $i \in [t]$.
        Let
        \[
          \tau \leftarrow \mathsf{ASig}(\mathbf{vk}, \mathbf{w}, m, \{(vk_i,\sigma_i)\}_{i=1}^{t})
          \quad\text{and}\quad
          \mathsf{avk} \leftarrow \mathsf{AKey}(\mathbf{vk}, \mathbf{w}).
        \]
        If the total stake of these $t$ parties meets the threshold, i.e.,
        \[
          \sum_{i=1}^{t} w_i \;\ge\; T,
        \]
        then
        \[
          \mathsf{AVer}(\mathsf{avk}, m, \tau, T) = 1.
        \]
\end{itemize}


% \begin{itemize}
%   \item \textbf{Aggregate key correctness.}
%         If $\mathsf{avk} \leftarrow \mathsf{AKey}(\mathbf{vk}, \mathbf{w})$, then
%         \[
%           \mathsf{ACheck}(\mathsf{avk}, \mathbf{vk}, \mathbf{w}) = 1.
%         \]

%   \item \textbf{Individual signature correctness.}
%         For every party index $i \in [N]$, let $(sk_i, vk_i)$ be generated by
%         $(sk_i, vk_i) \leftarrow \mathsf{KGen}(\mathsf{pp})$.
%         If $\sigma_i \leftarrow \mathsf{Sig}(vk_i, w_i, sk_i, j, m)$ and $\sigma_i \neq \bot$, then
%         \[
%           \mathsf{Ver}(vk_i, w_i, j, m, \sigma_i) = 1.
%         \]

%   \item \textbf{Aggregate signature correctness.}
%         Let $S \subseteq [N]$ be any subset of parties, and for each $i \in S$ let
%         $\sigma_i \leftarrow \mathsf{Sig}(vk_i, w_i, sk_i, j, m)$ be such that $\sigma_i \neq \bot$.
%         Let
%         \[
%           \tau \leftarrow \mathsf{ASig}(\mathbf{vk}, \mathbf{w}, m, \{(vk_i,\sigma_i)\}_{i \in S})
%           \quad\text{and}\quad
%           \mathsf{avk} \leftarrow \mathsf{AKey}(\mathbf{vk}, \mathbf{w}).
%         \]
%         If the total stake of $S$ meets the threshold, i.e.,
%         \[
%           \sum_{i \in S} w_i \;\ge\; T,
%         \]
%         then
%         \[
%           \mathsf{AVer}(\mathsf{avk}, m, \tau, T) = 1.
%         \]
% \end{itemize}


\subsubsection{Security}
\paragraph{Stake-based unforgeability $\alpha$-\textsf{UF-CMA}.}
In the unforgeability experiment, the challenger runs $\mathsf{Setup}$, generates honest keys $(sk_i,vk_i)$, fixes $\mathbf{w}$, computes $\mathsf{avk}$, and gives $(\mathsf{pp},\mathbf{vk},\mathbf{w},\mathsf{avk})$ to the adversary $\mathcal{A}$. The adversary can:
\begin{itemize}
  \item query \emph{signing oracles} for any honest user $i$, any $j$, and any message $m$, receiving $\mathsf{Sig}(vk_i,w_i,sk_i,j,m)$;
  \item \emph{corrupt} users to obtain $sk_i$ (and then act under those keys).
\end{itemize}
Let $C^\ast$ be the set of corrupted users. Let $H^\ast$ be the set of honest users that \emph{participated} on the forged topic (made at least one signing query for the challenge topic) under the forged aggregate key. The experiment \emph{charges} only successful eligibility (i.e., $\mathsf{Sig}\neq\bot$) and enforces the stake bound
\[
  \sum_{i\in C^\ast \cup H^\ast} w_i \;\le\; \alpha\, W .
\]
Finally, $\mathcal{A}$ outputs $(\mathbf{vk}^\ast,\mathsf{avk}^\ast,m^\ast,\tau^\ast)$; it wins if
\[
  \mathsf{ACheck}(\mathsf{avk}^\ast,\mathbf{vk}^\ast,\mathbf{w})=1,\quad
  \mathsf{AVer}(\mathsf{avk}^\ast,m^\ast,\tau^\ast,W)=1,
\]
and $\tau^\ast$ is \emph{fresh}, i.e., it cannot be obtained by aggregating local proofs returned (or derivable) from oracle responses. An STMS is $\alpha$-\textsf{UF-CMA} secure if every PPT adversary wins with only negligible probability.

\paragraph{Strengthened unforgeability $\alpha$-\textsf{UF-CMA-S}.}
As above, but the adversary may first register arbitrary public keys (and corresponding stake placement) before $\mathsf{AKey}$ is run (key-substitution/registration stage). The rest of the game is identical. Security requires negligible winning probability for all PPT adversaries.

% \medskip
% \noindent\emph{Notes on the model.}
% \begin{itemize}
%   \item Eligibility depends on the topic but not the body. This prevents grinding by changing $\textsf{body}$ while preserving eligibility.
%   \item We adopt the “do not charge unsuccessful attempts’’ rule: an ineligible honest user (output $\bot$) is \emph{not} counted toward the adversarial stake budget; this mirrors the bridge paper’s UF-CMA variant that balances probing versus corruption.
%   \item Completeness in aggregate correctness is enforced against one consistent set of indices; pairwise-distinct indices in $\mathsf{ASig}$ prevent double-counting a lottery win.
% \end{itemize}

% \subsubsection{Security}

% We follow the security definitions for stake-based threshold multi-signature schemes introduced in the bridge framework of Chaidos et al.\ and adapt them to the syntax of
% \[
%   \Pi_{\mathsf{STMS}}
%   =
%   (\mathsf{Setup}, \mathsf{KGen}, \mathsf{AKey}, \mathsf{ACheck},
%    \mathsf{Sig}, \mathsf{Ver}, \mathsf{ASig}, \mathsf{AVer}) .
% \]
% We assume that each message $m$ can be written as a pair $m = (\mathsf{topic}, \mathsf{body})$, where
% $\mathsf{topic}$ ranges over a comparatively small space and $\mathsf{body}$ ranges over a larger space.
% Eligibility to sign depends on the topic and on the aggregate verification key, but not on the body.

% Let $\alpha \in (0,1)$ be a bound on the adversarial fraction of the total stake.
% We capture security via two stake-bounded unforgeability notions under chosen-message attacks,
% denoted $\alpha$-UF-CMA and $\alpha$-UF-CMA-S, defined through a family of games.

% \paragraph{$\alpha$-UF-CMA game.}
% Fix a security parameter $\lambda \in \mathbb{N}$ and an adversary
% $\mathcal{A}$.
% The experiment
% \[
%   \mathsf{Exp}^{\alpha\text{-UF-CMA}}_{\Pi_{\mathsf{STMS}},\mathcal{A}}(\lambda)
% \]
% proceeds as follows.

% \begin{enumerate}
%   \item The challenger runs
%         $\mathsf{pp} \leftarrow \mathsf{Setup}(1^\lambda)$
%         and generates a vector of key pairs and stakes
%         \[
%           \bigl((sk_i,vk_i,w_i)\bigr)_{i \in [N]},
%         \]
%         for some $N \in \mathbb{N}$, by invoking
%         $(sk_i,vk_i) \leftarrow \mathsf{KGen}(\mathsf{pp})$ for each $i$ and assigning
%         positive stake values $w_i$.
%         Let $\mathbf{vk} = (vk_1,\dots,vk_N)$, $\mathbf{w} = (w_1,\dots,w_N)$,
%         and $W_{\mathsf{tot}} := \sum_{i=1}^{N} w_i$.
%         The challenger computes
%         \[
%           \mathsf{avk} \leftarrow \mathsf{AKey}(\mathbf{vk}, \mathbf{w})
%         \]
%         and gives $(\mathsf{pp},\mathbf{vk},\mathbf{w},\mathsf{avk})$ to $\mathcal{A}$.

%   \item The adversary $\mathcal{A}$ interacts adaptively with the following oracles:
%         \begin{itemize}
%           \item \emph{Corruption oracle}:
%                 on input an index $i \in [N]$, returns $sk_i$ and marks $i$ as corrupted.
%           \item \emph{Key-replacement oracle}:
%                 on input $(i, vk'_i)$, replaces $vk_i$ by $vk'_i$ and marks $i$ as controlled
%                 by $\mathcal{A}$.
%           \item \emph{Signing oracle}:
%                 on input $(i, m)$, where $m = (\mathsf{topic}, \mathsf{body})$,
%                 returns $\sigma \leftarrow \mathsf{Sig}(vk_i, w_i, sk_i, j, m)$
%                 for some lottery index $j$ chosen according to the scheme
%                 (or returns $\bot$ if the scheme declares $i$ ineligible for $(m,j)$).
%                 Whenever this oracle returns a non-$\bot$ signature for some message
%                 $m$ with topic $\mathsf{topic}$, the corresponding index $i$ is
%                 counted as under adversarial control for that topic.
%         \end{itemize}
%         We let $C_{\mathsf{topic}}$ denote the set of indices that are considered
%         under adversarial control for a topic $\mathsf{topic}$, i.e., indices that
%         have been corrupted, whose keys have been replaced, or for which the signing
%         oracle returned a non-$\bot$ signature on some message with that topic.

%   \item Eventually $\mathcal{A}$ outputs a forgery consisting of a message
%         $m^\ast = (\mathsf{topic}^\ast,\mathsf{body}^\ast)$,
%         an aggregate verification key $\mathsf{avk}^\ast$,
%         an aggregate signature $\tau^\ast$,
%         a stake threshold parameter $T^\ast$,
%         and a candidate description $(\mathbf{vk}^\ast,\mathbf{w}^\ast)$ of the
%         underlying verification keys and stakes.

%   \item The experiment checks the following conditions:
%         \begin{enumerate}
%           \item \emph{Well-formedness:}
%                 \[
%                   \mathsf{ACheck}(\mathsf{avk}^\ast,\mathbf{vk}^\ast,\mathbf{w}^\ast) = 1
%                   \quad\text{and}\quad
%                   \mathsf{AVer}(\mathsf{avk}^\ast, m^\ast, \tau^\ast, T^\ast) = 1.
%                 \]
%           \item \emph{Non-triviality:} the attack is not a trivial re-packaging of
%                 signatures already obtained from the oracles.  In particular, the
%                 aggregate signature $\tau^\ast$ cannot be obtained by merely aggregating
%                 signatures on $m^\ast$ that were returned by the signing oracle or
%                 produced using corrupted secret keys.
%           \item \emph{Stake bound:} writing
%                 $C^\ast := C_{\mathsf{topic}^\ast}$ for the set of indices counted as
%                 under adversarial control for topic $\mathsf{topic}^\ast$, and letting
%                 $W_{\mathsf{adv}} := \sum_{i \in C^\ast} w_i$, we require
%                 \[
%                   \frac{W_{\mathsf{adv}}}{W_{\mathsf{tot}}} \;\le\; \alpha .
%                 \]
%         \end{enumerate}
%         The experiment outputs $1$ if all checks pass, and $0$ otherwise.
% \end{enumerate}

% The \emph{$\alpha$-UF-CMA advantage} of $\mathcal{A}$ against $\Pi_{\mathsf{STMS}}$ is defined as
% \[
%   \mathsf{Adv}^{\alpha\text{-UF-CMA}}_{\Pi_{\mathsf{STMS}},\mathcal{A}}(\lambda)
%   :=
%   \Pr\bigl[
%     \mathsf{Exp}^{\alpha\text{-UF-CMA}}_{\Pi_{\mathsf{STMS}},\mathcal{A}}(\lambda) = 1
%   \bigr] .
% \]
% We say that $\Pi_{\mathsf{STMS}}$ is \emph{$\alpha$-UF-CMA secure} if this advantage is negligible
% in $\lambda$ for every probabilistic polynomial-time adversary $\mathcal{A}$.

% \paragraph{$\alpha$-UF-CMA-S game.}
% The static variant $\alpha$-UF-CMA-S is defined analogously, but with two modifications:
% (i) the adversary must perform all corruption and key-replacement queries before making any
% signing queries, and (ii) when computing the adversarial stake for the forgery, we only count
% stake belonging to indices that were corrupted or that were forced to sign the exact message
% $m^\ast$ being forged (rather than any message with the same topic).
% We denote the corresponding advantage by
% $\mathsf{Adv}^{\alpha\text{-UF-CMA-S}}_{\Pi_{\mathsf{STMS}},\mathcal{A}}(\lambda)$
% and say that $\Pi_{\mathsf{STMS}}$ is \emph{$\alpha$-UF-CMA-S secure} if this advantage is negligible
% for every probabilistic polynomial-time adversary.

% In what follows we assume that the underlying STMS scheme
% $\Pi_{\mathsf{STMS}}$ satisfies both $\alpha$-UF-CMA and $\alpha$-UF-CMA-S
% for suitable choices of the adversarial stake bound $\alpha$.


\subsection{Instantiation}

\subsubsection{Mithril}

We write the Mithril scheme as
\[
  \Pi_{\mathsf{Mithril}}
  =
  (\mathsf{Setup}, \mathsf{KGen}, \mathsf{AKey}, \mathsf{ACheck}, \mathsf{Sig}, \mathsf{ASig}, \mathsf{AVer}).
\]

\begin{itemize}

  \item $(\mathsf{pp}) \leftarrow \mathsf{Setup}(1^\lambda)$:\\
        run $\Pi_{\mathsf{BLS}}.\mathsf{Setup}(1^\lambda)$ to obtain $(q,G_1,G_2,G_T,g_1,g_2,e,H)$
        
        fix a dense mapping / eligibility function
        \[
        \mathsf{Eval} : ( msg, index, \sigma) \mapsto \{0,1\}^\lambda,
        \]
        % \[
        %   \mathsf{Eval} : G_1 \times \{1,\dots,\mu\} \times \{0,1\}^* \to \{0,1\}^\lambda,
        % \]
        a weighting function
        \[
          \varphi : \mathbb{N} \to [0,1],
        \]
        e.g. $\varphi(w) = 1-(1-f)^w$.

        Index range
        \[
          \mu \in \mathbb{N}, 
        \]

        a quorum or lottery parameter
        \[
          k \in \{1,\dots,\mu\}.
        \]
        Output
        \[
          \mathsf{pp}=(q,G_1,G_2,G_T,g_1,g_2,e,H,\mathsf{Eval},\varphi,\mu,k).
        \]

  \item $(sk,vk) \leftarrow \mathsf{KGen}(\mathsf{pp})$:\\
        parse $pp_{BLS}$ for BLS signature, run $\Pi_{\mathsf{BLS}}.\mathsf{KGen}(\mathsf{pp_{BLS}})$ 
        \[
          (sk,vk) \leftarrow \Pi_{\mathsf{BLS}}.\mathsf{KGen}(\mathsf{pp_{BLS}})
        \]
        Output $(sk,vk)$.

  \item $\mathsf{avk} \leftarrow \mathsf{AKey}(\mathbf{vk},\mathbf{w})$:\\
        given
        \[
          \mathbf{vk}=(vk_1,\dots,vk_N),
        \]
        \[
          \mathbf{w}=(w_1,\dots,w_N),
        \]
        set
        \[
          \mathsf{avk} \leftarrow \Pi_{\mathsf{MT}}.\mathsf{Create}(\mathbf{vk},\mathbf{w}).
        \]

  \item $b \leftarrow \mathsf{ACheck}(\mathsf{avk},\mathbf{vk},\mathbf{w})$, where $b\in\{0,1\}$:\\
        compute
        \[
          avk' \leftarrow \Pi_{\mathsf{MT}}.\mathsf{Create}(\mathbf{vk},\mathbf{w}).
        \]
        output b=1 iff avk=avk'. 

  \item $\pi \leftarrow \mathsf{Sig}(\mathsf{pp}, vk, w, sk, j, m)$:\\
        compute the BLS signature
        \[
          \sigma \leftarrow \Pi_{\mathsf{BLS}}.\mathsf{Sign}(\mathsf{pp}, sk, m),
        \]
        compute the eligibility value
        \[
          ev \leftarrow \mathsf{Eval}(\sigma, j, m),
        \]
        and check the threshold
        % \[
        %   ev \;<\; \bigl\lfloor \varphi(w)\cdot 2^\lambda \bigr\rfloor.
        % \]
        \[
          ev \;<\; \bigl\lfloor \varphi(w) \cdot 2^\lambda \bigr\rfloor.
        \]
        If true, output
        \[
          \pi=(j, \sigma, vk, w).
        \]
        Otherwise output $\bot$.




          \item $\tau \leftarrow \mathsf{ASig}(\mathsf{pp}, m, \{\pi_i\}_{i=1}^{t})$:\\
        parse each
        \[
          \pi_i=(j_i,\sigma_i,vk_i,w_i),
        \]
        with pairwise distinct
        \[
          j_i \in \{1,\dots,\mu\}.
        \]
        Set
        \[
          \sigma^\star \;=\; \prod_{i=1}^{t} \sigma_i \;\in\; G_1,
        \]
        \[
          vk^\star \;=\; \prod_{i=1}^{t} vk_i \;\in\; G_2.
        \]
        Derive a batching seed
        \[
          e_\sigma \;\leftarrow\; H\!\bigl(\sigma_1 \,\|\, \cdots \,\|\, \sigma_t \,\bigr),
        \]
        and coefficients for $i=1,\dots,t$
        \[
          \rho_i \;\leftarrow\; \mathsf{H_\lambda}(e_\sigma \,\|\, i) \in \mathbb{Z}_q^\ast.
        \]
        Compute the batch aggregates
        \[
          \mu \;=\; \prod_{i=1}^{t} \sigma_i^{\rho_i} \;\in\; G_1,
        \]
        \[
          ivk \;=\; \prod_{i=1}^{t} vk_i^{\rho_i} \;\in\; G_2.
        \]
        Output
        \[
          \tau \;=\; (\,\sigma^\star,\, vk^\star,\, \mu,\, ivk,\, e_\sigma).
        \]

  \item $b \leftarrow \mathsf{AVer}(\mathsf{pp}, \mathsf{avk}, m, \tau)$, where $b\in\{0,1\}$:\\
        parse
        \[
          \tau=(\sigma^\star, vk^\star, \mu, ivk, e_\sigma).
        \]
        Check that the indices are pairwise distinct.
        Compute
        \[
          h \leftarrow H(m).
        \]
        Verify the equation
        \[
          e(\sigma^\star, g_2) \;=\; e(h, vk^\star).
        \]
        Recompute the batching coefficients for $i=1,\dots,t$
        \[
          \rho_i \;\leftarrow\; \mathsf{H}(e_\sigma \,\|\, i) \in \mathbb{Z}_q^\ast,
        \]
        and the corresponding aggregated public key
        \[
          ivk' \;=\; \prod_{i=1}^{t} vk_i^{\rho_i} \;\in\; G_2.
        \]
        Check consistency of the submitted aggregated key
        \[
          ivk' \;=\; ivk.
        \]
        Verify the batched equation
        \[
          e(\mu, g_2) \;=\; e(h, ivk).
        \]
        Output
        \[
          b \;=\; 1 \text{ iff all checks above hold.}
        \]
\end{itemize}



% \section{}
% We write the Mithril scheme as
% \[
%   \Pi_{\mathsf{Mithril}}
%   =
%   (\mathsf{Setup}, \mathsf{KGen}, \mathsf{AKey}, \mathsf{ACheck},
%    \mathsf{Sig}, \mathsf{Ver}, \mathsf{ASig}, \mathsf{AVer}).
% \]

% We fix the following global parameters:
% \begin{itemize}
%   \item a security parameter $\lambda \in \mathbb{N}$;
%   \item a maximum number of participants $N \in \mathbb{N}$;
%   \item a stake vector
%         $\mathbf{w} = (w_1,\dots,w_N) \in \mathbb{N}^N$
%         with total stake $W := \sum_{i=1}^N w_i$;
%   \item an adversarial stake bound $\alpha \in (0,1/2)$.
% \end{itemize}

% \begin{itemize}

%   \item $(\mathsf{pp}) \leftarrow \mathsf{Setup}(1^\lambda)$: \\[0.2em]
%         Run $\Pi_{\mathsf{BLS}}.\mathsf{Setup}(1^\lambda)$ to obtain
%         \[
%           (q, G_1, G_2, G_T, g_1, g_2, e, H_{\mathsf{G}_1}),
%         \]
%         where $G_1,G_2,G_T$ are cyclic groups of order $q$,
%         $g_1 \in G_1$, $g_2 \in G_2$, $e : G_1 \times G_2 \to G_T$
%         is a bilinear pairing, and
%         $H_{\mathsf{G}_1} : \{0,1\}^* \to G_1$ is a hash-to-curve
%         function.

%         Fix a dense eligibility mapping
%         \[
%           Eval : G_1 \times \{1,\dots,\mu\} \times \{0,1\}^*
%                  \to \{0,1\}^\lambda,
%         \]
%         which we interpret as outputting a $\lambda$-bit string
%         (to be viewed as an integer in $\{0,\dots,2^\lambda-1\}$),
%         a weighting function
%         \[
%           \varphi : \mathbb{N} \to [0,1],
%         \]
%         e.g.\ $\varphi(w) = 1-(1-f)^w$ for some base rate $f$,

%         an index range
%         \[
%           \mu \in \mathbb{N},
%         \]
%         and a quorum / lottery parameter
%         \[
%           k \in \{1,\dots,\mu\}.
%         \]

%         Finally, fix an additional hash
%         \[
%           H_\lambda : \{0,1\}^* \to \mathbb{Z}_q^\ast
%         \]
%         for deriving batching coefficients.

%         Output
%         \[
%           \mathsf{pp} =
%           (q,G_1,G_2,G_T,g_1,g_2,e,H_{\mathsf{G}_1},Eval,\varphi,\mu,k,H_\lambda).
%         \]

%   \item $(sk,vk) \leftarrow \mathsf{KGen}(\mathsf{pp})$: \\[0.2em]
%         Run $\Pi_{\mathsf{BLS}}.\mathsf{KGen}(\mathsf{pp})$ to obtain
%         \[
%           sk \xleftarrow{\$} \mathbb{Z}_q, \qquad
%           vk = g_2^{sk} \in G_2.
%         \]
%         Output $(sk,vk)$.

%   \item $\mathsf{avk} \leftarrow \mathsf{AKey}(\mathbf{vk},\mathbf{w})$: \\[0.2em]
%         Given a vector of verification keys
%         \[
%           \mathbf{vk} = (vk_1,\dots,vk_N),
%         \]
%         and a stake vector
%         \[
%           \mathbf{w} = (w_1,\dots,w_N),
%         \]
%         compute the aggregate verification key using the Mithril
%         tree primitive:
%         \[
%           \mathsf{avk} \leftarrow
%           \Pi_{\mathsf{MT}}.\mathsf{Create}(\mathbf{vk},\mathbf{w}).
%         \]
%         Output $\mathsf{avk}$.

%   \item $b \leftarrow \mathsf{ACheck}(\mathsf{avk},\mathbf{vk},\mathbf{w})$,
%         where $b\in\{0,1\}$: \\[0.2em]
%         Recompute
%         \[
%           avk' \leftarrow \Pi_{\mathsf{MT}}.\mathsf{Create}(\mathbf{vk},\mathbf{w}).
%         \]
%         Output $b = 1$ iff $avk' = \mathsf{avk}$, and $b = 0$ otherwise.

%   \item $\pi \leftarrow \mathsf{Sig}(\mathsf{pp}, vk, w, sk, j, m)$: \\[0.2em]
%         Here $vk \in G_2$ is the signer's verification key,
%         $w \in \mathbb{N}$ is its local stake, $sk$ is its secret key,
%         $j \in \{1,\dots,\mu\}$ is a lottery index, and
%         $m \in \{0,1\}^*$ is the message.

%         Compute the BLS signature
%         \[
%           \sigma \leftarrow
%           \Pi_{\mathsf{BLS}}.\mathsf{Sign}(\mathsf{pp}, sk, m)
%           = H_{\mathsf{G}_1}(m)^{sk} \in G_1,
%         \]
%         compute the eligibility value
%         \[
%           ev \leftarrow Eval(\sigma, j, m) \in \{0,1\}^\lambda,
%         \]
%         and interpret $ev$ as an integer in $\{0,\dots,2^\lambda-1\}$.

%         Check the threshold
%         \[
%           ev \;<\; \bigl\lfloor \varphi(w) \cdot 2^\lambda \bigr\rfloor.
%         \]
%         If the inequality holds, output the local proof
%         \[
%           \pi = (j, \sigma, vk, w).
%         \]
%         Otherwise output $\bot$.

%   \item $b \leftarrow \mathsf{Ver}(\mathsf{pp}, vk, w, j, m, \pi)$,
%         where $b \in \{0,1\}$: \\[0.2em]
%         Parse
%         \[
%           \pi = (j', \sigma, vk', w').
%         \]
%         If $j' \neq j$ or $vk' \neq vk$ or $w' \neq w$, output $b=0$.
%         Otherwise, compute
%         \[
%           h \leftarrow H_{\mathsf{G}_1}(m)
%         \]
%         and verify the BLS signature
%         \[
%           e(\sigma, g_2) \stackrel{?}{=} e(h, vk).
%         \]
%         If this fails, output $b=0$.

%         Optionally recompute
%         \[
%           ev \leftarrow Eval(\sigma, j, m)
%         \]
%         and check the same threshold condition as in
%         $\mathsf{Sig}$; if it fails, output $b=0$.
%         Otherwise output $b=1$.

%   \item $\tau \leftarrow \mathsf{ASig}(\mathsf{pp}, m, \{\pi_i\}_{i=1}^{t})$: \\[0.2em]
%         Given a message $m \in \{0,1\}^*$ and a multiset of local
%         proofs
%         \[
%           \{\pi_i\}_{i=1}^{t},
%         \]
%         parse each
%         \[
%           \pi_i = (j_i,\sigma_i,vk_i,w_i),
%         \]
%         and require that the indices $j_i$ are pairwise distinct:
%         if $j_i = j_{i'}$ for some $i \neq i'$, output $\bot$.

%         For each $i=1,\dots,t$, optionally verify the local proof by
%         checking
%         \[
%           \mathsf{Ver}(\mathsf{pp}, vk_i, w_i, j_i, m, \pi_i) = 1;
%         \]
%         if any check fails, output $\bot$.

%         Set the (unweighted) aggregate signature and key
%         \[
%           \sigma^\star \;=\; \prod_{i=1}^{t} \sigma_i \;\in\; G_1,
%           \qquad
%           vk^\star \;=\; \prod_{i=1}^{t} vk_i \;\in\; G_2.
%         \]

%         Derive a batching seed
%         \[
%           e_\sigma \;\leftarrow\;
%           H(m \,\|\, \sigma_1 \,\|\, \cdots \,\|\, \sigma_t),
%         \]
%         and coefficients for $i=1,\dots,t$
%         \[
%           \rho_i \;\leftarrow\;
%           H_\lambda(e_\sigma \,\|\, i) \in \mathbb{Z}_q^\ast.
%         \]

%         Compute the batched aggregates
%         \[
%           \mu \;=\; \prod_{i=1}^{t} \sigma_i^{\rho_i} \;\in\; G_1,
%           \qquad
%           ivk \;=\; \prod_{i=1}^{t} vk_i^{\rho_i} \;\in\; G_2.
%         \]

%         Output
%         \[
%           \tau \;=\;
%           (\,\sigma^\star,\, vk^\star,\, \mu,\, ivk,\, e_\sigma\,).
%         \]

%   \item $b \leftarrow \mathsf{AVer}(\mathsf{pp}, \mathsf{avk}, m, \tau)$,
%         where $b\in\{0,1\}$: \\[0.2em]
%         Parse
%         \[
%           \tau = (\sigma^\star, vk^\star, \mu, ivk, e_\sigma).
%         \]

%         Compute
%         \[
%           h \leftarrow H_{\mathsf{G}_1}(m).
%         \]

%         First, verify the (unbatched) aggregate signature equation
%         \[
%           e(\sigma^\star, g_2) \stackrel{?}{=} e(h, vk^\star).
%         \]
%         If this fails, output $b=0$.

%         Recompute the batching coefficients (using public data)
%         for $i=1,\dots,t$
%         \[
%           \rho_i \;\leftarrow\; H_\lambda(e_\sigma \,\|\, i)
%           \in \mathbb{Z}_q^\ast,
%         \]
%         and the corresponding aggregated public key
%         \[
%           ivk' \;=\; \prod_{i=1}^{t} vk_i^{\rho_i} \;\in\; G_2.
%         \]
%         Check consistency of the submitted aggregated key
%         \[
%           ivk' \stackrel{?}{=} ivk.
%         \]
%         If this fails, output $b=0$.

%         Finally, verify the batched pairing equation
%         \[
%           e(\mu, g_2) \stackrel{?}{=} e(h, ivk).
%         \]
%         If this fails, output $b=0$.

%         Otherwise output $b = 1$.
% \end{itemize}






\end{document}
